% =====================================================================
% V-B. Ablación (a): etiquetas débiles  (E2a)
% Fuente: resultados/f6a_ablacion_resumen.csv
% =====================================================================

\subsection{Ablación (a): etiquetas débiles}
\label{sec:ablacion_debiles}

Con la pérdida combinada comparamos cuatro formas de usar los recortes débiles
(Tabla~\ref{tab:ablacion_a}). \emph{Solo dorados} es el modelo de la sección anterior.
\emph{Solo débiles} entrena únicamente con etiquetas de Otsu. \emph{Mezcla} añade 750 recortes
débiles por época a los 252 dorados. \emph{Preentreno + ajuste} entrena 30 épocas con débiles
(1\,000 por época) y luego ajusta solo con dorados, partiendo de esos pesos. Muestrear los débiles
en cada época mantiene un costo de minutos por corrida; a lo largo del entrenamiento cada
recorte débil se ve varias veces.

\begin{table}[t]
\centering
\caption{Ablación (a) en valid: media $\pm$ desviación sobre 3 semillas.}
\label{tab:ablacion_a}
\begin{tabular}{lccc}
\toprule
Estrategia & IoU & Precisión & Recall \\
\midrule
Otsu (referencia)        & 0.507                    & 0.634  & 0.716 \\
Solo débiles             & $0.562\pm0.009$          & 0.731 & 0.709 \\
Mezcla dorados + débiles & $0.571\pm0.004$          & 0.751 & 0.706 \\
Preentreno + ajuste      & $0.641\pm0.005$          & 0.823 & 0.744 \\
\textbf{Solo dorados}    & $\mathbf{0.644\pm0.002}$ & 0.822 & 0.749 \\
\bottomrule
\end{tabular}
\end{table}

Tres observaciones. Primera: la red entrenada \emph{solo} con etiquetas de Otsu supera a su
propio maestro (0.562 frente a 0.507). Al aprender un mapeo suave con contexto espacial
promedia parte del ruido de un umbral aplicado píxel a píxel. Segunda: mezclar débiles con
dorados \emph{empeora} el modelo ($-0.073$). Los débiles son tres veces más numerosos en cada
época y dominan el gradiente, de modo que la red hereda los errores sistemáticos de Otsu, en
especial la comisión en zonas áridas. Las cifras lo confirman: el recall apenas cambia entre
variantes (0.71--0.75), pero la comisión pasa de 0.178 con solo dorados a 0.249 con la mezcla y
a 0.269 con solo débiles. Tercera: preentrenar con débiles y ajustar con dorados
recupera el desempeño, pero queda empatado con usar solo dorados (la diferencia, $0.003$, es
menor que la desviación). Las etiquetas débiles no aportan información nueva: se derivan del
mismo SAR que ya ve la red, con una regla que la red supera. El modelo final usa, por tanto,
\textbf{solo los dorados con la pérdida combinada}.